% !TeX root = ./revletter_2.tex
\documentclass[12pt,a4paper]{article}

\usepackage[top=1.00in, bottom=1.0in, left=1.1in, right=1.1in]{geometry}
\usepackage{lineno}
\usepackage{graphicx}
\usepackage{natbib}
\usepackage{amsmath}
\usepackage{xr-hyper}
\externaldocument{paper_anon}
% \usepackage{hyperref}
\usepackage{gensymb}
\usepackage{todonotes}

% Cross referencing items
\newcommand{\R}[1]{\label{#1}\linelabel{#1}}
\newcommand{\lr}[1]{line~\ref{#1}}

\setlength\parindent{0pt}
\setlength {\marginparwidth }{2cm}
\begin{document}
%\bibliographystyle{}

% TO DO

% DISCUSS!

\emph{Reviewer and editor comments are in italics.} Our replies are in plain font.\\ 

% Senior editor comments start
{\bf Editor -- comments:} \\

\emph{Thank you for submitting your manuscript to Methods in Ecology and Evolution. I have now received the reviewers' reports and a recommendation from the Associate Editor who handled the review process. Copies of their reports are included below. This manuscript has the potential to make a valuable contribution to the area, but there are a number of significant concerns that need to be addressed. I have considered your paper in light of the comments received and I would like to invite you to prepare a major revision.}\\

We thank the editor for the opportunity to address the reviewer's concerns. Our point-by-point responses are below.\\
% Senior editor comments end

% Associate editor comments start
{\bf Associate Editor -- comments:} \\

\emph{The most important point is that the central claim of very high resolution can't currently be judged. Because no full-resolution stitched images are provided and the figures are low-resolution thumbnails without scale bars. Please include representative full-resolution examples for both a core and a cookie. Please also distinguish pixel sampling from effective optical resolution and reconcile the inconsistent values in the text.} \\

\todo{scan full resolution samples}
We agree that stitched images are necessary to evaluate the image quality. Full resolution scans for multiple species have been added in the supplementary material.\\

\emph{Two practical points need clarification: whether autofocusing only every 5 to 10 frames affects sharpness or stitching on uneven samples, and the generality of the particle-smearing step.}\\

Autofocusing every 5 to 10 frames is the result of the sample maintaining a nearly constant distance from the camera sensor. While it is possible for us to autofocus every frame, we have chosen to only refocus the lens if adjacent frames have significantly different focus scores. Capturing sequential frames without autofocusing is a indicates that the focus metric has not deviated from the past frame by more than 5\%, depending on configurable parameters. Therefore, the sharpness of the images are functionally identical when multiple frames are captured without refocusing the camera. We have clarified this in the manuscript between \lr{autofocusing_1} and \lr{autofocusing_2}.\\

\emph{Finally, please clarify how long and how intensively TIM has been in routine use, including approximate numbers of samples, species and operators, the period of use, and any failure modes. This would reassure the resource-limited labs you are targeting. Please also address the suggestion both reviewers list, and add scale references to the sample images throughout.} \\

We appreciate the associate editor for highlighting the need to provide more information about our usage of TIM. We have processed approximately 1,262 cores, 6 species, and 4 operators. Most of the scanning was done in a three month period where batches of 25 cores were scanned one to three times a day. Batches were often left scanning throughout the night without supervision. Additionally, we have provided the design files to another lab and they have built their own TIM. We have updated the manuscript to include these changes between \lr{usage_1} and \lr{usage_2}. \\

% Associate editor comments end

% Reviewer 1 comments start
{\bf Reviewer \# 1 -- comments:} \\

\emph{The most pressing concern is image quality. The pixel size of the Raspberry Pi HQ sensor (1.55 µm × 1.55 µm) is substantially smaller than sensors used in prosumer cameras currently used by existing systems. While the authors claim high resolution, no final stitched images are provided for evaluation to assess image quality. The low-resolution thumbnails in Figures 4 and 5, which also lack scale bars, make it impossible to assess whether the resolution gain includes reasonable image quality. This is particularly important given that the key innovation of such an advance in resolution is enabling anatomical measurements (moving forward current 4 microns from existing platforms is not necessary for annual rings or earlywood/latewood). I encourage the authors to include representative final images at full resolution.}\\

We thank the reviewer for highlighting that representative without full-resolution images, it is not possible to assess if the resolution gain improves the image quality appreciably. We have included multiple scans from different species in the supplementary material. \todo{Add scale bars to thumbnails} \todo{Add thumbnails with crop to indicate full resolution} Additional thumbnail images for the figures have been included to better showcase the resolution. \\ 
\todo{Scan full resolution samples}

\emph{Regarding Figure 6, the authors would benefit from revising the axes so that time appears as a function of sample size, separating cores from cookies, and increasing the number of samples to allow for variability assessment. In its current form the figure is difficult to interpret and I cannot get a good idea of the time necessary to run a core (still the most common way of dendrochronological sampling). Does a typical core 350 mm with 5 mm width(1750 mm2)  requires 150 minutes? I do not believe that if this were true you were presenting this article. So please, revise this table.} \\

\todo{Update figure with more samples, better axes, and separate cores from cookies}
We thank the reviewer for their insight to improve the legibility of the figure. We have changed the figure so that time appears as a function of the sample's size, separated cores from cookies, and increased the number of samples. \\

\emph{The particle smearing approach is imaginative but raises some concern. In my experience with similar systems, additional surface treatment has never been necessary for successful stitching. The example shown in Figure 5 corresponds to a gymnosperm with highly distinctive tracheid structures, I am shocked that stitching has problems with such an easy sample. It would be informative to know how TIM performs with diffuse-porous angiosperms such as Betula or Populus, where informative tracheid pattern is absent and wood is very white. Additionally, silicon carbide is a very hard material. In TIM configuration the machine is over the samples, so it should not be a problem, but I would be worried if a lab with optical material close to TIM is smearing sand in a routine fashion. I have lived in the coast for many years and there is nothing so disturbing than the way sand gets into every crack.}\\

We appreciate the reviewers comment and understand the concern that an easy sample would need additional surface treatment to be successfully stitched. We have changed the sample used in Figure 5 to be a harder to scan sample, more representative of a sample which would need surface treatment. Additionally, we have added specificity into the use case of the particle smearing method between \lr{particle_smearing_1} and \lr{particle_smearing_2}.\\

\emph{Regarding autofocusing, the authors indicate that autofocusing occurs approximately every 5-10 frames. Existing systems typically autofocus every frame. The authors should clarify whether this trade-off affects final image sharpness, particularly in samples with surface irregularities and how it affects to sharpness and stitching.}\\

We appreciate the concern of Reviewer 1 and have clarified how autofocusing every few frames affects the final image sharpness between \lr{autofocusing_1} and \lr{autofocusing_2}.\\

\emph{A minor but important point concerns the scale of measurement. The authors invest considerably (given the cost of TIM) in a precision micrometric scale for calibration, yet no scale reference appears in any of the sample images presented. For a measurement tool, every sample should include a scale reference. A simple millimetric reference would suffice (you can improve accuracy by measuring more milimeters in the scale) and that would allow readers to verify the reported DPI values independently.}\\

\todo{Scan ruler, add scale bars.}
We agree with the reviewer that a scale reference on the scan is necessary. Digital scale bars have been added to all samples. Additionally, a millimetric scale overlaid with a digital scale bar. Scanning a sample with a physical scale bar in the image is non-trivial due to the camera's narrow focal depth.\\

\emph{The statement in line 300 indicating a maximum core length of 1500 mm at full resolution warrants clarification, as this exceeds the practical dimensions of any standard increment core and it is complicated to work with a “cookie” that size.}\\

We thank the reviewer for pointing out that the maximum core length exceeds practical dimensions of increment borers. We have updated the text to clarify this on \lr{max_length}.\\

\emph{A final point, more philosophical in nature but worth raising. There is an important difference between a tool that works and a tool that is ready for others to use. In my experience, that difference only becomes apparent after months of routine use in a real lab setting: different samples, different operators, unexpected failures. The question for the authors is simple: is TIM currently part of their day-to-day data collection, or was this manuscript written closer to the assembly stage? This matters because the dendrochronology community has seen promising methods published before they were truly mature, and the labs that tried to adopt them paid the price in time and frustration. TIM is explicitly aimed at labs with limited economic capacity, which makes this question even more important, those labs have less margin to troubleshoot. The authors should provide an account of how long and how intensively the tool has been in us. This point is relevant to build confidence in the method.}\\

The reviewer raises valuable points when assessing the readiness for others to adopt TIM. While TIM does not have the polish of a commercial product, TIM has scanned over 1000 samples from 4 different operators over the past two years. Additionally, we shared our design files with another lab who have since recreated TIM and have been able to use it for their research. While TIM is not perfect, we have done our best to provide labs who want to replicate the design with the resources they need to reproduce TIM. All of TIM's design is open-source and we are willing to assist with reasonable inquiries regarding the replication process. These changes were made between \lr{usage_1} and \lr{usage_2}. \\

% Reviewer 1 comments end

% Reviewer 2 comments start 
{\bf Reviewer \# 2 -- comments:} \\

\emph{The system describes is very practical and I like that there are manual alternatives to the automated queuing approach. Likewise for cookies/disk material, specific radii can be manually selected to save time and file storage space/size. It would be useful to note what happens when a storage limit is hit -what does the system do in such a case? OK, I see that uncompressed array files are saved – but does this allow stitching?}\\

We thank the reviewer for highlighting the need for clarification regarding file size limits. When the file size limit is reached, the system still stitches the sample in the same method but it is saved as a memory mapped NumPy array instead of a tiff image. The manuscript has been updated with this clarification on \lr{large_stitch}.\\

\emph{Its great that the study even tests the use of a grain enhancer like silicon carbide. I presume that more dramatic stains would not affect imaging and stitching contrast settings?}\\

More dramatic stains would not affect the imaging and stitching contrast settings. We chose to use silicon carbide as it was more reversable than staining.\\

\emph{It’s a little unclear with the current notation what “21 140 DPI” means – should it be 21,140 DPI? Also, one version of your abstract indicates 0.83 pixels per micron and the other 1.2 microns per pixel. Similarly for Point 3, is the value \$3,000 USD? The problem seems to be a space has been introduced to replace the thousand separators: L54 – 6 400 DPI rather than 6,400 DPI. See L67, L69, L88, L89… etc. throughout the manuscript.}\\

We changed the notation to use commas in our number notation. Additionally the abstract was revised to only display microns per pixel. \todo{Update the submission abstract to include this} \\

\emph{
    \begin{itemize}
        \item L96: Change “current” to currently 
        \item L129: Should read “…and have a planar surface for…”
        \item L137: Small point admittedly, but conventionally the reference to a figure should appear before the figure itself occurs. Same for Figure 3 which appears before the reference to it.
        \item L160: Gu et al. should not be within parenthesis given the current structure, only the publication date (2015).
        \item L185: Clarification needed. “…same threshold was used across all cores for 5 species…” Is the point that the sequence worked for variation caused by having different species, or that the transition from one species to another was handled?
    \end{itemize}}

The manuscript has been updated to address these line items. \\
    
\begin{itemize}
        \item \lr{line_change_1}: Changed to "currently"
        \item \lr{line_change_2}: Changed to "...have a planar surface..."
        \item Moved figure reference above figure \todo{figure out how to do this}
        \item \lr{line_change_3}: Changed citation.
        \item \lr{line_change_4}: Added clarification that the sequence works when transition between species in the same batch.
    \end{itemize}


\emph{Section 2.4.2: This description needs a little rephrasing for better clarity. Quite a number of sentences repeat words (code was generated to convert code, package similarly. Again final sentence -code was generated for codebase integration and was noted in code… etc. I think I understand in general what you mean, but I’m not sure. Please rephrase more carefully.}\\

We thank the reviewer for addressing the repetition of words in the manuscript. We have rephrased the section to improve clarity between \lr{llm_rewording_1} and \lr{llm_rewording_2}.

\emph{Figure S1: The scales in the two images are very difficult to read for the ordinary human, slightly larger font would be a help. The scale in the (b) image currently seems to be out of sync with the image.} \\

We have updated the figures to make font larger and resync the images. \todo{Update figure with scale bar, increase font, etc.}\\


\newpage
% \bibliography{..//..//..//..//vitisbiblio}
% \bibliographystyle{apa}

\end{document}